\markdownRendererDocumentBegin
\markdownRendererSectionBegin
\markdownRendererHeadingOne{Arboles}\markdownRendererInterblockSeparator
{}\markdownRendererInputFencedCode{./_markdown_main/96c255a98b50da304ad9cf0250fa4913.verbatim}{cpp}{cpp}\markdownRendererParagraphSeparator
{}\markdownRendererSectionBegin
\markdownRendererHeadingTwo{Binary Tree Traversals}\markdownRendererInterblockSeparator
{}Recorridos clásicos de árbol: preorden, inorden y postorden.\markdownRendererInterblockSeparator
{}\markdownRendererInputFencedCode{./_markdown_main/b6097f481e5893045f315cca89c4e40d.verbatim}{cpp}{cpp}\markdownRendererInterblockSeparator
{}\markdownRendererDollarSign{}\markdownRendererBackslash{}clearpage\markdownRendererDollarSign{}\markdownRendererInterblockSeparator
{}
\markdownRendererSectionEnd 
\markdownRendererSectionEnd \markdownRendererSectionBegin
\markdownRendererHeadingOne{Calculating Tree Diameter}\markdownRendererInterblockSeparator
{}Usa dos DFS para obtener el diámetro del árbol (camino más largo entre dos nodos) en O(n).\markdownRendererInterblockSeparator
{}\markdownRendererInputFencedCode{./_markdown_main/373d4cb1dbb80ba1b3787bdd26beac26.verbatim}{cpp}{cpp}\markdownRendererInterblockSeparator
{}\markdownRendererDollarSign{}\markdownRendererBackslash{}clearpage\markdownRendererDollarSign{}\markdownRendererInterblockSeparator
{}\markdownRendererSectionBegin
\markdownRendererHeadingTwo{All longest paths}\markdownRendererInterblockSeparator
{}Calcula para cada nodo la distancia máxima a cualquier otro nodo (reroot DP clásico). Complejidad O(n).\markdownRendererInterblockSeparator
{}\markdownRendererInputFencedCode{./_markdown_main/2f864b49212ab2aef0e473ee45004afe.verbatim}{cpp}{cpp}\markdownRendererInterblockSeparator
{}\markdownRendererDollarSign{}\markdownRendererBackslash{}clearpage\markdownRendererDollarSign{}\markdownRendererInterblockSeparator
{}
\markdownRendererSectionEnd \markdownRendererSectionBegin
\markdownRendererHeadingTwo{Lowest Common Ancestor (LCA)}\markdownRendererInterblockSeparator
{}Encuentra el ancestro común más bajo entre dos nodos a y b usando binary lifting.\markdownRendererInterblockSeparator
{}\markdownRendererInputFencedCode{./_markdown_main/f17a83fb5157579d41e348d5e11c7cf0.verbatim}{cpp}{cpp}\markdownRendererInterblockSeparator
{}
\markdownRendererSectionEnd \markdownRendererSectionBegin
\markdownRendererHeadingTwo{Balanced Factor}\markdownRendererInterblockSeparator
{}Calcula el factor de balance (altura izquierda - altura derecha) para verificar si el árbol está equilibrado.\markdownRendererSoftLineBreak
{}(No usa rotaciones, solo cálculo conceptual.)\markdownRendererInterblockSeparator
{}\markdownRendererInputFencedCode{./_markdown_main/631aa8fde7c6d18f4cecc8e4ab001f37.verbatim}{cpp}{cpp}
\markdownRendererSectionEnd 
\markdownRendererSectionEnd \markdownRendererDocumentEnd